% =====================================================================
% TERRA -- modelo de relatorio de analise
%
% Compilar de dentro de docs/report/:
%   latexmk relatorio.tex
% O latexmk vive em /Library/TeX/texbin, fora do PATH padrao. O .latexmkrc
% deste diretorio encontra terra.sty em ../user-manual/tex-files.
%
% COMO USAR. Copie este arquivo para o relatorio novo e troque cada
% \preencher{...} e \preencherbloco{...} pelo dado. Toda lacuna aparece no
% PDF como caixa creme entre colchetes e e listada no log. Para emitir,
% troque a opcao `rascunho' por `final': com ela uma lacuna restante
% interrompe a compilacao, e o PDF liberado nao chega a existir incompleto.
%
% DE ONDE VEM CADA NUMERO. As lacunas de resultado nomeiam o campo da
% execucao que as alimenta (overall_pct, peak_water_fraction_pct, ...), como
% ele esta em internal/analysis/types.go e frontend/src/lib/types.ts. Um
% numero do relatorio e copiado da execucao, nunca recalculado a mao, e o
% identificador da execucao vai na secao de reprodutibilidade.
%
% SECAO DE PRODUTO NAO EXECUTADO SAI INTEIRA. Um relatorio de cobertura do
% solo nao leva uma secao de agua vazia "para manter a estrutura".
%
% A ORDEM DAS SECOES e a sequencia do leitor de escrita.md (secao 4): isto
% e para mim (abertura, sintese), o que isto faz (pergunta e escopo), eu
% acredito no numero (metodo, resultados, limitacoes), como reproduzo
% (reprodutibilidade, apendices).
% =====================================================================

\documentclass[11pt,a4paper]{article}

\usepackage[utf8]{inputenc}
\usepackage[T1]{fontenc}
\usepackage[brazilian]{babel}
% Opcoes: english | final | rascunho. Ver o cabecalho de terra-report.sty.
\usepackage[rascunho]{terra-report}

\graphicspath{{figuras/}}

% --- identificacao (ISO 7200) ----------------------------------------
% Numeracao sugerida: TERRA-RA-<ano>-<sequencial de tres digitos>. A revisao
% comeca em A e avanca a cada reemissao; o historico no fim registra o que
% mudou em cada uma.
\titulo{Cobertura do solo e situação socioambiental de \preencher{nome da área}}
\titulocurto{\preencher{nome da área}}
\subtitulo{\preencher{município}, \preencher{UF}. Safra \preencher{safra},
  Sentinel-2 de \preencher{início} a \preencher{fim}}
\identificacao{TERRA-RA-\preencher{ano}-\preencher{nnn}}
\revisao{A}
\emissao{\preencher{data de emissão}}
\situacao{\emelaboracao}
\titular{\preencher{instituição ou empresa}}
\destinatario{\preencher{quem recebe}}
\responsavel{\preencher{autor}}
\aprovacao{\preencher{quem aprova}}

\begin{document}

\terraabertura

% --- sintese ---------------------------------------------------------
% Tres a cinco achados. Cada um com o numero, a unidade e a condicao em que
% foi medido. Nenhum qualificador sem numero ("expressiva", "baixa"): o
% leitor que so ler esta caixa sai com os valores, e nao com uma impressao.
\secaolivre{Síntese}

Este relatório descreve \preencher{produtos executados} sobre
\preencher{nome da área} (\preencher{área} ha, \preencher{n} talhões), a partir
de \preencher{n\_dates} cenas Sentinel-2 L2A entre \preencher{início} e
\preencher{fim}.

\begin{sintese}
  \item \preencherbloco{achado 1: valor, unidade e condição de medida.
        Ex.: classe dominante e sua fração, com a concordância com o MapBiomas}
  \item \preencherbloco{achado 2}
  \item \preencherbloco{achado 3}
\end{sintese}

% Ficha de procedencia, na primeira pagina pelo motivo que terra.sty da:
% "com que dado?" e a pergunta que decide se vale ler o resto.
\begin{ficha}
  \item[Área] \preencher{nome}, \preencher{area\_ha} ha;
        \preencher{n} talhões, limites de \preencher{delineação ou arquivo importado}
  \item[Imagens] Sentinel-2 L2A, catálogo STAC do Planetary Computer;
        \preencher{start} a \preencher{end}, nuvem abaixo de
        \preencher{max\_cloud}\,\%, \preencher{uma cena por mês ou todas}
  \item[Referência] MapBiomas Coleção \preencher{n}, ano \preencher{year}
  \item[Produtos] \preencher{classificação, fenologia, saúde, água,
        sobreposição, mineral}
  \item[Software] TERRA \preencher{versão}; execuções listadas na
        seção~\ref{sec:repro}
\end{ficha}

\newpage

% =====================================================================
\section{Pergunta e escopo}
% A pergunta do destinatario, em uma ou duas frases, e o que o relatorio
% deixa de fora. Nao descreva a tela: diga o que a analise responde.

\preencherbloco{A pergunta que o destinatário fez, e a decisão que depende
dela. Ex.: quanto da área está sob cultura anual na safra, e se algum
talhão sobrepõe cadastro público de restrição.}

\preencherbloco{O que o relatório não cobre. Ex.: produtividade, sanidade
por patógeno, regularidade fundiária.}

% Estas duas ressalvas vem do proprio codigo (sidecar/terra/overlap e
% frontend/src/lib/health.ts). Mantenha a que corresponder a um produto
% executado e apague a outra.
\begin{aviso}
Sobreposição com cadastro público não constitui constatação de
irregularidade, e vale para a data em que os cadastros foram lidos.
Desvio de índice de vegetação não constitui diagnóstico: semeadura mais
tardia, outra cultura ou pousio produzem o mesmo desvio que estresse.
\end{aviso}

% =====================================================================
\section{Área e talhões}

\begin{figure}[htbp]\centering
  \figurapendente[8cm]{Mapa da área com os limites dos talhões, identificados
  pelo nome usado nas tabelas. Escala gráfica, norte e sistema de referência}
  \legenda{\preencher{O que o mapa sustenta}}{Limites de
  \preencher{origem}. Sistema de referência \preencher{EPSG}.
  Fundo: composição de cor verdadeira de \preencher{data}.}
  \label{fig:area}
\end{figure}

\begin{table}[htbp]\centering\small
  \caption{Talhões da área (\preencher{n} talhões,
  \preencher{área total} ha).}\label{tab:talhoes}
  \begin{tabular}{@{}lS[table-format=4.1]ll@{}}
    \arrayrulecolor{terratrim}\toprule
    \cab{Talhão} & {\cab{Área (ha)}} & \cab{Cultura declarada} & \cab{Origem do limite} \\
    \midrule
    {\preencher{nome}} & {\preencher{ha}} & {\preencher{cultura}} & {\preencher{origem}} \\
    \bottomrule
  \end{tabular}
\end{table}

% =====================================================================
\section{Dados e método}
% Uma subsecao por produto EXECUTADO. O texto de cada uma afirma o que o
% sidecar faz; o arquivo que o sustenta esta no comentario ao lado. Se o
% codigo mudar, o texto muda junto.

\subsection{Aquisição}
% frontend/src/lib/methodBrief.ts, acquisition(). Criterio de selecao, conforme
% monthly_best: "Uma cena por mês, a de menor cobertura de nuvem" ou "Todas
% as cenas abaixo do limite".
Cenas Sentinel-2 L2A obtidas pelo catálogo STAC do Planetary Computer, de
\preencher{start} a \preencher{end}, com cobertura de nuvem abaixo de
\preencher{max\_cloud}\,\%. \preencher{critério de seleção}. Apenas a janela do polígono e as
bandas usadas pelo modelo são lidas. \preencher{n\_dates} datas atenderam ao
filtro (lista no apêndice~\ref{ap:cenas}).

\subsection{Cobertura do solo}
% sidecar/terra/landcover/mapbiomas.py
Classificação por pixel com \preencher{modelo}, na legenda de cinco classes
MapBiomas (3, 21, 25, 39, 41). A concordância com o mapa MapBiomas
\citep{souza2020} é medida por exatidão global, exatidão do produtor e do
usuário, com intervalo de Wilson a 95\,\% \citep{brown2001}, e pela divisão da
discordância em quantidade e alocação \citep{pontius2011}. O coeficiente kappa
não é calculado, pelo argumento de \citet{pontius2011}.

\subsection{Fenologia}
% sidecar/terra/phenology.py
Série temporal de NDVI médio por talhão, suavizada por Savitzky--Golay
\citep{savitzky1964}, da qual se extraem início (SOS), pico (POS) e fim (EOS)
da estação, em dia do ano, e a duração (LOS), em dias.

\subsection{Saúde da vegetação}
% sidecar/terra/health/series.py; frontend/src/lib/health.ts
NDVI e NDRE de cada data da estação comparados a duas referências: as
estações anteriores (\preencher{baseline\_years}), nas datas a até
\preencher{window\_days} dias do mesmo dia do ano, em desvios-padrão; e os
demais talhões do conjunto numa data comum, contra a mediana, em
desvios-padrão robustos ($1{,}4826 \times$ desvio absoluto mediano;
\citealp{rousseeuw1993}). Os cortes de leitura são 1 e 2 desvios-padrão.

\subsection{Água superficial}
% sidecar/terra/water/indices.py
Índice \preencher{NDWI | MNDWI} \citep{xu2006} por data, com limiar
\preencher{Otsu | fixo} \citep{otsu1979}. A ocorrência é a fração das datas
válidas em que o pixel foi classificado como água; água efêmera é a presente
em parte das datas, e persistente a presente na maioria.

\subsection{Sobreposição com cadastros públicos}
% sidecar/terra/overlap; frontend/src/lib/overlap.ts
Interseção da área com PRODES, DETER, embargos do IBAMA e do ICMBio, terras
indígenas, unidades de conservação e CAR, como estavam em
\preencher{read\_at}. Feições sobrepostas de um mesmo cadastro contam uma vez.
Cadastro que não pôde ser lido é reportado como não lido, e não como zero.

\subsection{Mapa mineral}
% sidecar/terra/mineral
Identificação de minerais de superfície pelo sistema especialista Tetracorder
\citep{clark2003} sobre reflectância EMIT, em células de 60 m, com uma resposta
por grupo Tetracorder. Seleção de passagem por NDVI e índice de absorção de
celulose.

% =====================================================================
\section{Resultados}

\subsection{Cobertura do solo}

\begin{figure}[htbp]\centering
  \figurapendente[8cm]{Mapa de classes e mapa de confiança lado a lado,
  painéis a e b}
  \legenda{\preencher{Achado que o mapa sustenta}}{\textbf{a},~classe
  prevista por pixel; \textbf{b},~confiança do classificador, de 0 a 1.
  Cores das classes conforme a legenda MapBiomas.}
  \label{fig:classes}
\end{figure}

% Area por classe. Numeros copiados de class_stats (area_ha, pct).
\begin{table}[htbp]\centering\small
  \caption{Área por classe prevista (\preencher{n\_pixels} pixels de
  \preencher{pixel\_size\_m} m).}\label{tab:classes}
  \begin{tabular}{@{}lS[table-format=5.1]S[table-format=3.1]@{}}
    \arrayrulecolor{terratrim}\toprule
    \cab{Classe} & {\cab{Área (ha)}} & {\cab{Fração (\%)}} \\
    \midrule
    {\preencher{name}} & {\preencher{area\_ha}} & {\preencher{pct}} \\
    \midrule
    Total & {\preencher{ha}} & 100,0 \\
    \bottomrule
  \end{tabular}
\end{table}

% Concordancia. Numeros de agreement: overall_pct, overall_ci,
% quantity_disagreement_pct, allocation_disagreement_pct, per_class.
A exatidão global contra o MapBiomas \preencher{year} é de
\preencher{overall\_pct}\,\% (IC 95\,\%: \preencher{overall\_ci}), sobre
\preencher{n\_reference\_cells} células de referência. A discordância divide-se
em \preencher{quantity\_disagreement\_pct}\,\% de quantidade e
\preencher{allocation\_disagreement\_pct}\,\% de alocação.

\begin{table}[htbp]\centering\small
  \caption{Exatidão por classe contra o MapBiomas, com intervalo de Wilson a
  95\,\%.}\label{tab:exatidao}
  \begin{tabular}{@{}lllrr@{}}
    \arrayrulecolor{terratrim}\toprule
    \cab{Classe} & \cab{Produtor (\%)} & \cab{Usuário (\%)} &
    \cab{$n$ ref.} & \cab{$n$ prev.} \\
    \midrule
    \preencher{name} & \preencher{producers\_pct [CI]} &
    \preencher{users\_pct [CI]} & \preencher{n} & \preencher{n} \\
    \bottomrule
  \end{tabular}
\end{table}

\begin{objecao}
A comparação mede concordância com outro mapa, não exatidão contra
verificação de campo. Não há desenho amostral por trás dela; a avaliação de
exatidão com amostragem é a de \citet{olofsson2014}, e esta não é uma.
\end{objecao}

\subsection{Fenologia}

\begin{figure}[htbp]\centering
  \figurapendente[6cm]{Série de NDVI por talhão ao longo da estação, com SOS,
  POS e EOS marcados}
  \legenda{\preencher{Achado que a série sustenta}}{NDVI médio por data,
  $n$ = \preencher{n\_dates} datas. Linha: série suavizada; marcadores:
  médias observadas; faixa: \preencher{desvio-padrão}.}
  \label{fig:fenologia}
\end{figure}

\begin{aviso}
NDVI marca emergência e senescência, não semeadura e colheita. A semeadura
antecede a emergência e a colheita sucede a senescência; com revisita de 5 dias
e nuvem, a incerteza das datas é de \preencher{n} dias nesta série.
\end{aviso}

\subsection{Saúde da vegetação}

\preencherbloco{Talhões com desvio abaixo de $-1$ e de $-2$ desvios-padrão na
data \preencher{map\_date}, com o valor de cada um e a referência usada.}

\begin{figure}[htbp]\centering
  \figurapendente[7cm]{Mapa de anomalia de NDVI na data mais recente com
  referência}
  \legenda{\preencher{Achado}}{NDVI menos a mediana das estações anteriores na
  mesma época do ano, por célula, em \preencher{map\_date}. Escala divergente
  centrada em zero, de $-$\preencher{map\_range} a $+$\preencher{map\_range}.}
  \label{fig:saude}
\end{figure}

\subsection{Água superficial}

A fração máxima de água foi de \preencher{peak\_water\_fraction\_pct}\,\% da
área, em \preencher{peak\_date}. Água persistente ocupou
\preencher{persistent\_area\_ha} ha, e efêmera \preencher{ephemeral\_area\_ha}
ha, em \preencher{n\_dates} datas.

\subsection{Sobreposição com cadastros públicos}

% Registros em REGISTER_ORDER (frontend/src/lib/overlap.ts); publicador
% conforme o campo `publisher' de cada camada. Cadastro nao
% lido: escreva "não lido" na coluna de area, nunca 0.
\begin{table}[htbp]\centering\small
  \caption{Área sobreposta a cada cadastro, lida em \preencher{read\_at}.}
  \label{tab:cadastros}
  \begin{tabularx}{\linewidth}{@{}L{4.2cm}L{3.8cm}rrY@{}}
    \arrayrulecolor{terratrim}\toprule
    \cab{Cadastro} & \cab{Publicador} & \cab{Área (ha)} & \cab{Feições} &
    \cab{Nota} \\
    \midrule
    PRODES              & INPE, TerraBrasilis & \preencher{ha} & \preencher{n} & \\
    DETER               & INPE, TerraBrasilis & \preencher{ha} & \preencher{n} & \\
    Embargo IBAMA       & IBAMA   & \preencher{ha} & \preencher{n} & \\
    Embargo ICMBio      & ICMBio, publicado pelo IBAMA & \preencher{ha} & \preencher{n} & \\
    Terra indígena      & FUNAI   & \preencher{ha} & \preencher{n} & \\
    Unidade de conservação & MMA, CNUC & \preencher{ha} & \preencher{n} & \\
    CAR                 & SFB, SICAR & \preencher{ha} & \preencher{n} & \\
    \bottomrule
  \end{tabularx}
\end{table}

Supressão PRODES após o marco do Código Florestal (\preencher{forest\_code\_cutoff}):
\preencher{ha} ha. Após o marco do regulamento europeu EUDR
(\preencher{eudr\_cutoff}): \preencher{ha} ha.

\subsection{Mapa mineral}

\preencherbloco{Grupos Tetracorder identificados, fração da área com resposta
em cada um, e a concordância com o produto EMIT L2B da mesma passagem.}

% =====================================================================
\section{Resultados por talhão}\label{sec:talhoes}
% As colunas do Field table do compositor (frontend/src/lib/fieldTable.ts).
% Celula sem dado fica com travessao, e a nota da tabela diz por que.

\begin{xltabular}{\linewidth}{@{}L{2.2cm}YL{1.35cm}L{1.35cm}L{1.35cm}rrr@{}}
  \caption{Resultados por talhão. SOS, POS e EOS em data; confiança média do
  classificador, de 0 a 1; desvio de NDVI em desvios-padrão contra as
  estações anteriores.}\label{tab:porTalhao}\\
  \arrayrulecolor{terratrim}\toprule
  \cab{Talhão} & \cab{Classe dominante} & \cab{SOS} & \cab{POS} & \cab{EOS} &
  \cab{Conf.} & \cab{$z$ NDVI} & \cab{Cadastro (ha)} \\
  \midrule\endfirsthead
  \toprule
  \cab{Talhão} & \cab{Classe dominante} & \cab{SOS} & \cab{POS} & \cab{EOS} &
  \cab{Conf.} & \cab{$z$ NDVI} & \cab{Cadastro (ha)} \\
  \midrule\endhead
  \bottomrule\endfoot
  \preencher{nome} & \preencher{classe (\%)} & \preencher{data} &
  \preencher{data} & \preencher{data} & \preencher{0,00} &
  \preencher{z} & \preencher{ha} \\
\end{xltabular}

% =====================================================================
\section{Limitações}
% Cada limitacao com a medida e com a direcao em que move o resultado.
% Apague as que nao se aplicam; acrescente as que a execucao reportou.
\begin{limitacoes}
  \item[Resolução espacial] Pixel de 10 m (0,01 ha). Talhões com borda
        irregular perdem \preencher{n}\,\% da área em pixels mistos de borda.
  \item[Nuvem] \preencher{n\_dates} datas válidas no período; o maior
        intervalo sem observação foi de \preencher{n} dias, entre
        \preencher{data} e \preencher{data}.
  \item[Confiança do classificador] A confiança é a probabilidade atribuída
        pelo modelo à classe prevista. Uma confiança de 0,67 não corresponde a
        67\,\% de acerto.
  \item[Referência] MapBiomas ano \preencher{year} contra imagens de
        \preencher{ano}. Mudança de uso entre os dois anos aparece como
        discordância.
  \item[Validade] Os cadastros públicos refletem \preencher{read\_at}; um
        embargo ou alerta posterior não está neste relatório.
  \item[Campo] \preencher{Nenhuma verificação de campo | n pontos
        verificados em data}.
\end{limitacoes}

% =====================================================================
\section{Reprodutibilidade}\label{sec:repro}

\begin{ficha}
  \item[TERRA] versão \preencher{AppVersion}, commit \preencher{hash}
  \item[Ambiente] \preencher{requirements.lock.txt, hash}
  \item[Projeto] \preencher{arquivo .terra}
  \item[Execuções] \preencher{run\_id} (\preencher{tipo}),
        \preencher{data e hora UTC}
\end{ficha}

% O pedido que o TERRA enviou ao sidecar, copiado da execucao. O ambiente
% `parametros' cita o texto como esta; lacuna nao funciona dentro dele, por
% isso a lacuna fica fora e o exemplo comentado mostra a forma:
%
% \begin{parametros}
% {"start": "2025-09-01", "end": "2026-04-30", "max_cloud": 30,
%  "monthly_best": true, "model_kind": "...", "area_id": "..."}
% \end{parametros}
\preencherbloco{Pedido JSON de cada execução, no ambiente parametros.}

% =====================================================================
\appendix

\section{Matriz de confusão}\label{ap:matriz}
% agreement.matrix, com agreement.matrix_classes como rotulos. Linhas:
% prevista; colunas: referencia (a ordem de mapbiomas.py).
\preencherbloco{Matriz de confusão em número de células, linhas = classe
prevista, colunas = classe de referência.}

\section{Cenas utilizadas}\label{ap:cenas}
% DataCubeScene: id, date, cloud_cover, tile, satellite.
\begin{xltabular}{\linewidth}{@{}YL{2.2cm}rL{1.6cm}L{1.8cm}@{}}
  \arrayrulecolor{terratrim}\toprule
  \cab{Identificador} & \cab{Data} & \cab{Nuvem (\%)} & \cab{Tile} &
  \cab{Satélite} \\
  \midrule\endfirsthead
  \toprule
  \cab{Identificador} & \cab{Data} & \cab{Nuvem (\%)} & \cab{Tile} &
  \cab{Satélite} \\
  \midrule\endhead
  \bottomrule\endfoot
  \preencher{id} & \preencher{date} & \preencher{cloud\_cover} &
  \preencher{tile} & \preencher{satellite} \\
\end{xltabular}

% =====================================================================
\begin{thebibliography}{99}
\small

\bibitem[Brown et~al.(2001)]{brown2001}
Brown, L.~D., Cai, T.~T., DasGupta, A. (2001). Interval estimation for a
binomial proportion. \textit{Statistical Science}, 16(2), 101--133.

\bibitem[Clark et~al.(2003)]{clark2003}
Clark, R.~N., Swayze, G.~A., Livo, K.~E., et~al. (2003). Imaging spectroscopy:
Earth and planetary remote sensing with the USGS Tetracorder and expert
systems. \textit{Journal of Geophysical Research: Planets}, 108(E12), 5131.

\bibitem[Olofsson et~al.(2014)]{olofsson2014}
Olofsson, P., Foody, G.~M., Herold, M., Stehman, S.~V., Woodcock, C.~E.,
Wulder, M.~A. (2014). Good practices for estimating area and assessing
accuracy of land change. \textit{Remote Sensing of Environment}, 148, 42--57.

\bibitem[Otsu(1979)]{otsu1979}
Otsu, N. (1979). A threshold selection method from gray-level histograms.
\textit{IEEE Transactions on Systems, Man, and Cybernetics}, 9(1), 62--66.

\bibitem[Pontius e Millones(2011)]{pontius2011}
Pontius, R.~G., Millones, M. (2011). Death to Kappa: birth of quantity
disagreement and allocation disagreement for accuracy assessment.
\textit{International Journal of Remote Sensing}, 32(15), 4407--4429.

\bibitem[Rousseeuw e Croux(1993)]{rousseeuw1993}
Rousseeuw, P.~J., Croux, C. (1993). Alternatives to the median absolute
deviation. \textit{Journal of the American Statistical Association}, 88(424),
1273--1283.

\bibitem[Savitzky e Golay(1964)]{savitzky1964}
Savitzky, A., Golay, M.~J.~E. (1964). Smoothing and differentiation of data by
simplified least squares procedures. \textit{Analytical Chemistry}, 36(8),
1627--1639.

\bibitem[Souza et~al.(2020)]{souza2020}
Souza, C.~M., Shimbo, J.~Z., Rosa, M.~R., et~al. (2020). Reconstructing three
decades of land use and land cover changes in Brazilian biomes with Landsat
archive and Earth Engine. \textit{Remote Sensing}, 12(17), 2735.

\bibitem[Xu(2006)]{xu2006}
Xu, H. (2006). Modification of normalised difference water index (NDWI) to
enhance open water features in remotely sensed imagery.
\textit{International Journal of Remote Sensing}, 27(14), 3025--3033.

\end{thebibliography}

% =====================================================================
\secaolivre{Controle do documento}

\begin{historico}
  \revisaolinha{A}{\preencher{data}}{Emissão inicial.}{\preencher{autor}}
\end{historico}

\assinaturas{\preencher{nome}}{\preencher{função}}
            {\preencher{nome}}{\preencher{função}}

\end{document}
